\chapter{Thermodynamics with and without Dynamics}

\small
\begin{quote}
The general study of the dynamical stability of fluid motions \dots\ is concerned
with the approach to equilibrium, not with the equilibrium itself. (after S.\ Chandrasekhar)
\end{quote}
\begin{quote}
A theory is the more impressive the greater the simplicity of its premises, the
more different kinds of things it relates, and the more extended its area of
applicability. Therefore the deep impression that classical thermodynamics made
upon me. It is the only physical theory of universal content which I am convinced
will never be overthrown. (Albert Einstein)
\end{quote}
\normalsize

\section{A Theory of Endpoints}

Classical thermodynamics, in the form handed down from Clausius, Kelvin and Gibbs,
is a theory of \emph{states} and not of \emph{processes}. This is easy to overlook,
because the textbooks are full of arrows and cycles and talk of ``processes''
throughout. But look closely at what the equations contain: there is no time
variable in them. The fundamental relation
\begin{equation}
  dU = T\,dS - p\,dV
  \label{eq:fund}
\end{equation}
is a relation between the \emph{differentials of state functions}, not a law of
evolution. It tells you how the equilibrium quantities of one resting state differ
from those of a neighbouring resting state. It does not tell you how, or how fast,
or along which path in real space and time, the one is reached from the other.

The ``processes'' of classical thermodynamics are \emph{quasi-static}: idealised as
an infinitely slow succession of equilibrium states, so slow that at every instant
the system is, to the accuracy of the theory, already at rest. A quasi-static
process is a curve drawn on the equilibrium surface, not a trajectory through it in
time. There is no rate in it, because a rate would require a clock, and the theory
carries none. The Second Law in this setting appears as an \emph{inequality between
states},
\begin{equation}
  \Delta S \ge 0,
  \label{eq:second}
\end{equation}
which marks a direction on the equilibrium surface but says nothing about the speed
of travel or the route taken.

\section{How the Dynamics is Made to Disappear}

Equilibrium statistical mechanics, which was built to give Clausius's entropy a
mechanical meaning, removes the dynamics even more deliberately. The partition
function
\begin{equation}
  Z = \sum_i e^{-\beta E_i}
  \label{eq:partition}
\end{equation}
contains no time whatsoever: it is a weighted count of states. The \emph{ergodic
hypothesis} is then invoked precisely so that the time average along the true
mechanical trajectory --- the thing one would have to compute by integrating the
equations of motion --- may be replaced by an average over this static ensemble of
states. The entire device exists in order to obtain equilibrium numbers
\emph{without} solving the dynamics. That is the source of its power, and it is
exactly the source of its silence on everything we found to be essential in the
preceding chapters: the approach, the selection of the realised state, the rate of
dissipation.

Where, then, does the dynamics live in the standard subject? It lives in
\emph{separate} theories, bolted on from outside and usually as empirical closures:
Fourier's law of heat conduction, Fick's law of diffusion, the Navier--Stokes and
Euler equations, the Boltzmann transport equation, the near-equilibrium reciprocal
relations of Onsager. These carry the rates and the trajectories that equilibrium
thermodynamics lacks, but they sit \emph{outside} the equilibrium core and connect
to it cleanly only in the linear regime of small departures from rest. The standard
edifice is therefore predominantly a \emph{dynamics-free theory of endpoints}, with
the motion either idealised away (quasi-static), averaged away (ergodicity), or
appended afterwards as phenomenological transport.

\section{The Approach, not the State, is What Matters}

Why should this trouble us? Because, as we have seen throughout this book, the
physically essential content of a thermodynamic problem is the \emph{approach to
equilibrium}, not the equilibrium state in isolation.

The reason is simple and decisive. The conservation laws and constraints of a given
problem fix only a \emph{family} of admissible end-states --- a whole manifold of
configurations compatible with the prescribed energy, volume and composition.
Equilibrium thermodynamics, by itself, hands you this family; it does not tell you
\emph{which} member of it a given initial condition will reach. That selection is
made entirely by the \emph{dynamics}: by the path, the boundary conditions, the rate
of driving, and above all by the dissipation along the way. Quote an equilibrium
state without the process that produces it and you have named a member of a family
without saying which one --- an under-determined object.

This is not a philosopher's quibble. It is the everyday situation of metastability,
hysteresis and the glassy state, where the realised state is fixed by the approach
and is emphatically \emph{not} the global minimiser of the free energy: supercooled
water that does not freeze, diamond that does not relax to graphite, the remanent
magnetisation of iron, the glass that never finds its crystal. In every such case
the equilibrium-minimisation answer is simply \emph{wrong} about what is observed,
and only an account of the dynamics gets it right. And the Second Law itself ---
the irreversible production of entropy, or in the language of this book the
turbulent dissipation --- is a statement about the \emph{approach}. At the endpoint
there is no production; the endpoint is precisely where production has ceased.

It is therefore close to meaningless to speak of the equilibrium state without also
speaking of how it is reached. The one legitimate exception is the degenerate case
in which the equilibrium is \emph{unique} and \emph{rapidly} reached --- a dilute
gas in a box, a simple chemical equilibrium with a single deep minimum --- so that
every path arrives at the same place and the history may be forgotten. This
history-independence is the genuine power of state functions, and it is the whole
domain in which equilibrium thermodynamics earns its keep as a standalone theory.
But notice what the exception concedes: the endpoint description is complete only
when there is nothing left to choose. The moment the selection among end-states is
genuinely non-trivial --- competing minima, slow relaxation, long-range forces such
as self-gravitation --- the equilibrium description is incomplete without the
dynamics, and the interesting physics has migrated entirely into the approach.

\section{The Right Way Up}

The ordering of the standard subject is thus inverted. It takes the dynamics-free
theory of endpoints as fundamental and treats the approach --- where the Second Law
actually operates and where the realised state is selected --- as an afterthought to
be supplied by separate transport laws. The fact that so much can be said
\emph{without} dynamics is, on reflection, not a sign of depth but a measure of how
much has been idealised away.

The programme of this book turns the picture the right way up. We make the
\emph{dynamics} primary: the compressible Euler equations solved by Direct Finite
Element Simulation, with turbulent dissipation arising unavoidably from
least-squares stabilisation as the finite-precision residue of an unresolvable fine
scale. The Second Law then comes out not as a postulate about entropy but as a
\emph{theorem about the computation}: energy cascades from resolved large scales to
unresolved small scales, and finite precision forbids the reverse, so the
dissipation is one-way. Equilibrium thermodynamics is recovered from this as the
\emph{long-time limit} --- the stationary endpoint at which the resolvable free
energy is exhausted and all macroscopic fluxes have died out. The state functions
are then bookkeeping of that endpoint, valid as a self-contained theory exactly in
the lucky case where the choice of endpoint collapses to one.

In short: the general theory is the theory of the \emph{approach}; equilibrium is
what remains to be said when there is nothing left to choose. A thermodynamics that
begins with the dynamics and obtains equilibrium as its asymptote is complete where
the classical subject is silent, and it pays for this completeness in the only
honest currency available --- the explicit, computable, irreversible dissipation of
the real motion.
